\section*{अध्याय \ref{ch:b2:orbital-forcing} --- कक्षीय प्रणोदन और हिमयुग}

\begin{solution}{exo:b2:orbital-forcing:1}
(क) $1.0167/0.9833 = 1.0340$: अपसौर उपसौर से \qty{3.4}{\percent} अधिक दूर है।
(ख) $a(1 + e) - a(1 - e) = 2ea = 0.0334\,a$, यानी \qty{3.34}{\percent}, $a$ के सापेक्ष, जो प्रायः उद्धृत 3.4 प्रतिशत
के निकट है। (ग) सूर्य-प्रकाश दूरी के व्युत्क्रम वर्ग के अनुसार घटता है: $1.0340^2 = 1.069$,
उपसौर पर लगभग \qty{7}{\percent} अधिक सूर्य-प्रकाश।
\end{solution}

\begin{solution}{exo:b2:orbital-forcing:2}
(क) आज $0.017236 \times \sin\ang{101.37} = 0.017236 \times 0.9804 = 0.01690$; \num{10000} वर्ष पहले
$0.019971 \times \sin\ang{293.86} = 0.019971 \times (-0.9145) = -0.01826$; \num{6000} वर्ष पहले
$0.019249 \times \sin\ang{359.99} \approx 0$, जो तालिकाओं के मान हैं। (ख) उपसौर पर सूर्य का
देशांतर $\varpi + \ang{180}$ है: आज \ang{281.37}, \ang{473.86} $-$ \ang{360} $=$ \ang{113.86}
\num{10000} वर्ष पहले, \ang{179.99} \num{6000} वर्ष पहले। (ग) \num{10000} वर्ष पहले उपसौर
जून अयनांत के \ang{24} बाद पड़ता था, तीनों में उसके सबसे निकट, उत्तरी ग्रीष्म में;
\num{6000} वर्ष पहले वह सितंबर विषुव पर पड़ता था, जहाँ \omterm{def:b2:orbital-forcing:elements}{जलवायवीय पुरस्सरण} शून्य है।
\end{solution}

\begin{solution}{exo:b2:orbital-forcing:3}
(क) आवृत्तियाँ जुड़ती हैं: $1/25.77 + 1/112 = 0.03880 + 0.00893 = \qty{0.04773}{kyr^{-1}}$, यानी
$1/0.04773 = 21.0$ हज़ार वर्ष का आवर्तकाल। (ख) परिणाम मुख्य आवर्तकालों के परास, 19.0 से 23.7 हज़ार
वर्ष, के भीतर पड़ता है। एक अकेली संख्या उन्हें पुनरुत्पन्न नहीं कर सकती: कक्षा का घूमना
स्वयं भिन्न आवृत्तियों वाली अनेक गतियों का योग है, जिनमें से हर एक अक्ष के पुरस्सरण के साथ
मिलकर \omterm{def:b2:orbital-forcing:elements}{जलवायवीय पुरस्सरण} का एक आवर्तकाल देती है; औसत प्रायः उद्धृत
\num{23000} वर्ष है।
\end{solution}

\begin{solution}{exo:b2:orbital-forcing:4}
$\sqrt{1 - 0.0167^2} = 0.99986$। \ang{22.1} के लिए: $1360 \times 0.37622/(\pi \times 0.99986) =
\qty{162.9}{W.m^{-2}}$; \ang{24.5} के लिए: $1360 \times 0.41469/(\pi \times 0.99986) =
\qty{179.5}{W.m^{-2}}$। अनुपात $\sin\ang{24.5}/\sin\ang{22.1} = 1.102$ है: दोनों चरम मानों के बीच ध्रुव पर
वार्षिक माध्य \qty{10.2}{\percent}, लगभग \qty{16.7}{W.m^{-2}}, बढ़ता है।
\end{solution}

\begin{solution}{exo:b2:orbital-forcing:5}
(क) जून अयनांत पर $(a/r)^2 = \bigl((1 - e\sin\varpi)/(1 - e^2)\bigr)^2$ है। आज, जब $e\sin\varpi
= 0.016898$ है, यह गुणक $(0.983102/0.999703)^2 = 0.9671$ है, जबकि $1 - 0.033796 = 0.9662$;
\num{116000} वर्ष पहले $e\sin\varpi = 0.043939$ और $(0.956061/0.998065)^2 = 0.9176$, जबकि $1 -
0.087878 = 0.9121$। (ख) दिसंबर अयनांत पर $(a/r)^2 = \bigl((1 + e\sin\varpi)/(1 -
e^2)\bigr)^2$: आज $1.0347$, जबकि $1.0338$; \num{116000} वर्ष पहले $1.0940$, जबकि $1.0879$।
(ग) उपेक्षित पद $e^2$ की कोटि के हैं: $(1 - x)^2/(1 - e^2)^2 \approx 1 - 2x + x^2 + 2e^2$,
जहाँ $x = e\sin\varpi$, यहाँ लगभग $3e^2$। वे आज $3 \times 0.0172^2 = 0.0009$ और $3
\times 0.0440^2 = 0.0058$ \num{116000} वर्ष पहले हैं, जो पाए गए अंतरों, $0.0009$ और
$0.0055$, के निकट हैं: जब उत्केंद्रता ढाई गुना
बड़ी हो तो सन्निकटन छह गुना अधिक खराब होता है।
\end{solution}

\begin{solution}{exo:b2:orbital-forcing:6}
जून अयनांत पर $\delta = \varepsilon = \ang{23.446}$। $\cos h_0 = -\tan\ang{65}\tan\ang{23.446}
= -2.14451 \times 0.43369 = -0.93006$, अतः $h_0 = \ang{158.44} = \qty{2.7654}{rad}$। दिन
$2 \times 158.44/15 = 21.13$~घंटे चलता है (पृथ्वी प्रति घंटे \ang{15} घूमती है)। $\sin\ang{65} = 0.90631$,
$\cos\ang{65} = 0.42262$, $\sin\delta = 0.39788$, $\cos\delta = 0.91744$ और $\sin h_0 = 0.36742$ के साथ,
कोष्ठक $2.7654 \times 0.90631 \times 0.39788 + 0.42262 \times 0.91744 \times 0.36742 = 0.99720
+ 0.14246 = 1.13966$ है। माध्य दूरी पर $\bar Q = (1360/\pi) \times 1.13966 = 432.90 \times 1.13966
= \qty{493.4}{W.m^{-2}}$; \cref{exo:b2:orbital-forcing:5} के दूरी गुणक के साथ, $493.4 \times
0.9671 = \qty{477.1}{W.m^{-2}}$। जाँच के रूप में, \ang{60} उ पर वही चरण \qty{475.96}{W.m^{-2}} देते हैं, जहाँ तालिकाएँ \num{475.95} देती हैं।
\end{solution}

\begin{solution}{exo:b2:orbital-forcing:7}
(क) मार्च विषुव $\lambda = 0$ है, सितंबर विषुव $\lambda = \pi$। $\sqrt{1 - e^2}
\approx 1$ के साथ, अवधि $\frac{T}{2\pi}\int_0^{\pi}\bigl(1 + 2e\cos(\lambda -
\varpi)\bigr)\,\mathrm d\lambda = \frac{T}{2\pi}\bigl(\pi + 2e[\sin(\lambda - \varpi)]_0^{\pi}\bigr)
= \frac{T}{2\pi}(\pi + 4e\sin\varpi) = \frac T2 + \frac{2T}{\pi}e\sin\varpi$ है। (ख) दूसरा अर्ध-वर्ष
उतना ही छोटा है, इसलिए अंतर $(4T/\pi)\,e\sin\varpi = (4 \times
365.2422/\pi) \times 0.016898 = 465.04 \times 0.016898 = 7.86$~दिन है। (ग) वसंत और ग्रीष्म मिलकर
शरद और शीत ऋतु से $4.5 + 3 = 7.5$~दिन अधिक हैं; उद्धृत दोनों अंतर पूर्णांकित हैं, और
सहमति उस पूर्णांकन के भीतर है।
\end{solution}

\begin{solution}{exo:b2:orbital-forcing:8}
(क) $e = 0.0046$ के लिए: $\sqrt{1 - e^2} = 0.999989$ और $340/0.999989 = \qty{340.004}{W.m^{-2}}$।
$e = 0.0565$ के लिए: $\sqrt{1 - 0.003192} = 0.998403$ और $340/0.998403 = \qty{340.544}{W.m^{-2}}$। (ख)
अंतर \qty{0.54}{W.m^{-2}} है, \qty{0.16}{\percent}। (ग) \ang{65} उ पर मध्य-जुलाई सूर्यातप
$497 - 388 = \qty{109}{W.m^{-2}}$ झूला, अपने माध्य का लगभग \qty{25}{\percent} और वैश्विक वार्षिक माध्य के
परिवर्तन का दो सौ गुना: कक्षा जलवायु पर सूर्य-प्रकाश को ऋतुओं और अक्षांशों के बीच
खिसकाकर कार्य करती है, योग बदलकर नहीं।
\end{solution}

\begin{solution}{exo:b2:orbital-forcing:9}
$\omega = 2\pi/P$ और $\tau = \qty{10}{kyr}$ के साथ: $P = 23$ के लिए $\omega\tau = 2.732$, सापेक्ष
आयाम $1/\sqrt{1 + 7.464} = 0.344$ और पश्चता $\arctan(2.732)/\omega = 1.2199/0.2732 = 4.47$~\unit{kyr};
$P = 41$ के लिए $\omega\tau = 1.532$, आयाम $0.546$, पश्चता $0.9926/0.1532 = 6.48$~\unit{kyr}; $P =
100$ के लिए $\omega\tau = 0.628$, आयाम $0.847$, पश्चता $0.5610/0.0628 = 8.93$~\unit{kyr}। हिमचादर
धीमे प्रणोदन को वरीयता देती है: वह \qty{85}{\percent} तक \num{100000}-वर्षीय चक्र का अनुसरण करती है पर
पुरस्सरण-चक्र के केवल एक-तिहाई का, क्योंकि तेज़ प्रणोदन के उलटने से पहले उसके पास बढ़ने या सिकुड़ने के लिए
बहुत कम समय होता है। वह एक निम्न-पारक फ़िल्टर की तरह कार्य करती है।
\end{solution}

\begin{solution}{exo:b2:orbital-forcing:10}
(क) समीकरण रैखिक है: यदि $V_1$, $A_1\sin\omega_1t$ का और $V_2$,
$A_2\sin\omega_2t$ का उत्तर है, तो $V_1 + V_2$ उनके योग का उत्तर है। \cref{prop:b2:orbital-forcing:lag} से $V_1$ और
$V_2$ आवृत्तियों $\omega_1$ और $\omega_2$ के ज्यावक्र हैं (जब क्षणिक भाग मिट जाए), इसलिए
अनुक्रिया में कोई अन्य आवृत्ति नहीं होती। (ख) $V \approx B_1\sin(\omega_1t - \phi_1) +
B_2\sin(\omega_2t - \phi_2)$ के साथ,
\[
\begin{aligned}
  V^2 &= \tfrac12(B_1^2 + B_2^2) - \tfrac12B_1^2\cos(2\omega_1t - 2\phi_1) -
  \tfrac12B_2^2\cos(2\omega_2t - 2\phi_2)\\
  &\quad + B_1B_2\bigl[\cos\bigl((\omega_1 - \omega_2)t - \phi_1 + \phi_2\bigr) -
  \cos\bigl((\omega_1 + \omega_2)t - \phi_1 - \phi_2\bigr)\bigr],
\end{aligned}
\]
अतः पद $kV^2$ समीकरण को अंतर और योग आवृत्तियों पर, और दोगुनी
आवृत्तियों पर, प्रणोदित करता है। (ग) $1/(1/19.0 - 1/23.7) = 1/(0.05263 - 0.04219) = 95.8$ हज़ार वर्ष; 19.0
और 22.4 का युग्म 125 हज़ार वर्ष देता है। पुरस्सरण पदों की अंतर आवृत्तियाँ उत्केंद्रता की पट्टी में
पड़ती हैं, जैसा कि उन्हें पड़ना ही चाहिए, क्योंकि \omterm{def:b2:orbital-forcing:elements}{जलवायवीय पुरस्सरण} $e\sin\varpi$ है। इसलिए एक अरैखिक
हिमचादर पुरस्सरण से \num{100000}-वर्षीय लय निकाल सकती है; एक रैखिक हिमचादर
नहीं। \num{100000}-वर्षीय चक्र अरैखिक व्यवहार के साक्ष्य हैं, यद्यपि केवल यह परिकलन
यह नहीं दिखाता कि कौन-सा।
\end{solution}

\begin{solution}{exo:b2:orbital-forcing:11}
(क) \cref{prop:b2:orbital-forcing:annual} की उपपत्ति की तरह, $(a/r)^2\,\mathrm dt =
T\,\mathrm d\lambda/(2\pi\sqrt{1 - e^2})$, और $\sin\delta = \sin\varepsilon\sin\lambda$
$0 < \lambda < \pi$ के लिए धनात्मक है। अतः $E = \frac{S_0T\sin\varepsilon}{2\pi\sqrt{1 - e^2}}\int_0^{\pi}\sin\lambda
\,\mathrm d\lambda = S_0T\sin\varepsilon/(\pi\sqrt{1 - e^2})$। (ख) $T = 365.25 \times 86400 =
\qty{3.1558e7}{s}$, $S_0T/\pi = \qty{1.3661e10}{J.m^{-2}}$; \ang{22.1} के लिए $E = 1.3661 \times 10^{10}
\times 0.37622/0.99986 = \qty{5.14}{GJ.m^{-2}}$; \ang{24.5} के लिए $\qty{5.67}{GJ.m^{-2}}$,
\qty{10}{\percent} अधिक। (ग) पुरस्सरण तय करता है कि वही ऊर्जा समय में कैसे फैलती है: उपसौर पर पड़ने वाली
ग्रीष्म छोटी और अधिक तीव्र होती है, अपसौर पर पड़ने वाली ग्रीष्म लंबी और दुर्बल। पिघलना
केवल \qty{0}{\celsius} से ऊपर शुरू होता है, इसलिए एक छोटी तीव्र ग्रीष्म उतना ही योग पाने वाली
एक लंबी ठंडी ग्रीष्म से अधिक पिघला सकती है।
\end{solution}

\begin{solution}{exo:b2:orbital-forcing:12}
(क) $3.93\log_2(279.4/189.7) = 3.93\log_2 1.4729 = 3.93 \times 0.5586 = \qty{2.20}{W.m^{-2}}$। (ख)
सबसे छोटी वृद्धि 194 से \qty{266}{ppm} तक है: $3.93\log_2 1.3711 = 3.93 \times 0.4554 =
\qty{1.79}{W.m^{-2}}$; सबसे बड़ी 188 से \qty{282}{ppm} तक: $3.93\log_2 1.5 = 3.93 \times 0.5850 =
\qty{2.30}{W.m^{-2}}$। (ग) कार्बन डाइऑक्साइड का प्रणोदन वैश्विक वार्षिक सूर्यातप के सबसे बड़े
परिवर्तन, \qty{0.54}{W.m^{-2}}, का चार गुना है, और वह हर जगह और पूरे वर्ष कार्य करता है; \ang{65} उ पर ग्रीष्म
सूर्यातप लगभग \qty{109}{W.m^{-2}} झूलता है, पर एक अक्षांश पर और एक ऋतु में।
कक्षा ट्रिगर है, उत्तरी हिमचादरों की ग्रीष्म ऋतुओं के माध्यम से; कार्बन डाइऑक्साइड,
\omterm{def:b2:precambrian-earth:ice-albedo}{हिम--एल्बिडो प्रतिपुष्टि} के साथ, एक प्रवर्धक है जो परिवर्तन को पूरे ग्लोब पर फैलाता है।
\end{solution}

\begin{solution}{pb:b2:orbital-forcing:1}
\textbf{1.} तब $0.0440 \times \sin\ang{92.7} = 0.0440 \times 0.99889 = 0.04395$; अब $0.0172 \times
\sin\ang{101.4} = 0.0172 \times 0.98027 = 0.01686$, लगभग 2.6 गुना कम।
\textbf{2.} तब $1 - 0.0440 = 0.956$ और $1 + 0.0440 = 1.044$; अब $0.9828$ और $1.0172$।
\textbf{3.} $\lambda = \varpi + \ang{180}$ पर: तब \ang{272.7}, अब \ang{281.4}, दोनों
दिसंबर अयनांत ($\lambda = \ang{270}$) के निकट; तब पृथ्वी उसके तीन दिन से भी कम बाद उपसौर से
गुज़री, और जून अयनांत पर अपसौर से।
\textbf{4.} तब $(1.044/0.956)^2 = 1.0921^2 = 1.193$, अब $(1.0172/0.9828)^2 = 1.0350^2 = 1.071$:
तब उपसौर पर \qty{19}{\percent} अधिक सूर्य-प्रकाश, अब \qty{7}{\percent}।
\textbf{5.} कम, \ang{22.52} बनाम \ang{23.45}: उच्च अक्षांशों पर ग्रीष्म का सूर्य नीचा था,
और उच्च अक्षांशों को ग्रीष्म में और पूरे वर्ष में कम सूर्य-प्रकाश मिला।
\textbf{6.} तब $340/\sqrt{1 - 0.001936} = 340/0.999032 = \qty{340.33}{W.m^{-2}}$; अब $340/0.999852
= \qty{340.05}{W.m^{-2}}$। अंतर, \qty{0.28}{W.m^{-2}} यानी प्रति हज़ार एक भाग से कम,
नगण्य है।
\textbf{7.} $\lambda = \ang{90}$ पर, $\cos(\ang{90} - \varpi) = \sin\varpi$, यह
\cref{prop:b2:orbital-forcing:distance} में रखिए।
\textbf{8.} तब $(1 - 0.001936)/(1 - 0.04395) = 0.998064/0.956049 = 1.04395$; अब $(1 -
0.000296)/(1 - 0.01686) = 0.999704/0.983139 = 1.01685$।
\textbf{9.} तब $1/1.04395^2 = 0.91758$; अब $1/1.01685^2 = 0.96714$।
\textbf{10.} तब $1 - 2 \times 0.04395 = 0.91210$; अब $1 - 2 \times 0.01686 = 0.96628$:
सन्निकटन तब $0.0055$ से चूकता है, अब $0.0009$ से।
\textbf{11.} तब $1 - 0.91758 = \qty{8.2}{\percent}$ दुर्बल, अब \qty{3.3}{\percent}।
\textbf{12.} अयनांत पर $\delta = \varepsilon$: तब \ang{22.52}, अब \ang{23.45}।
\textbf{13.} तब $\cos h_0 = -2.14451 \times 0.41462 = -0.88916$, $h_0 = \ang{152.77} =
\qty{2.6663}{rad}$; अब $\cos h_0 = -2.14451 \times 0.43378 = -0.93023$, $h_0 = \ang{158.47} =
\qty{2.7659}{rad}$।
\textbf{14.} $2h_0/\ang{15}$ प्रति घंटा: तब $20.37$~घंटे, अब $21.13$~घंटे।
\textbf{15.} तब $2.6663 \times 0.90631 \times 0.38301 + 0.42262 \times 0.92375 \times 0.45759 =
0.92554 + 0.17864 = 1.10418$ (अपूर्णांकित \num{1.10417}), और $432.90 \times 1.10417 =
\qty{478.0}{W.m^{-2}}$; अब $2.7659 \times 0.90631 \times 0.39795 + 0.42262 \times 0.91741 \times
0.36697 = 0.99757 + 0.14228 = 1.13985$ (अपूर्णांकित \num{1.13982}), और $432.90 \times 1.13982 = \qty{493.4}{W.m^{-2}}$।
\textbf{16.} तब $478.0 \times 0.91758 = \qty{438.6}{W.m^{-2}}$; अब $493.4 \times 0.96714 =
\qty{477.2}{W.m^{-2}}$, तब से \qty{38.6}{W.m^{-2}} अधिक।
\textbf{17.} मध्य-जुलाई में सूर्य का देशांतर \ang{120} है: $\sin\delta = \sin\varepsilon\sin\ang{120}$,
यानी लगभग \ang{20} की क्रांति, \ang{23} के बजाय, इसलिए छोटा दिन और नीचा सूर्य: \ang{65} उ पर
दैनिक-माध्य सूर्यातप अयनांत के निकट सबसे अधिक होता है, जब क्रांति सबसे बड़ी होती है।
\textbf{18.} आज की दूरी के साथ, केवल तिर्यकता का परिवर्तन $(493.4 - 478.0) \times
0.96714 = \qty{14.9}{W.m^{-2}}$ घटाता है; पुरानी तिर्यकता के साथ, केवल दूरी $478.0 \times
(0.96714 - 0.91758) = \qty{23.7}{W.m^{-2}}$ घटाती है। दोनों का योग \qty{38.6}{W.m^{-2}} है; दूरी का
योगदान लगभग तीन-पाँचवाँ रहा।
\textbf{19.} आकस्मिक परिवर्तन के बाद हिमचादर उसका $1 - \mathrm e^{-t/\tau}$ प्राप्त कर चुकी होती है:
$1 - \mathrm e^{-0.5} = 0.39$ \qty{5}{kyr} के बाद, $1 - \mathrm e^{-1} = 0.63$ \qty{10}{kyr} के बाद,
$1 - \mathrm e^{-2} = 0.86$ \qty{20}{kyr} के बाद।
\textbf{20.} $\omega\tau = 2\pi \times 10/23 = 2.732$: हिमचादर प्रणोदन के $1/\sqrt{1 + 2.732^2} =
0.34$ का अनुसरण करती है, $\arctan(2.732)/\omega = 4.5$~\unit{kyr} देर से।
\textbf{21.} आधा पुरस्सरण-चक्र, लगभग 11 से \qty{12}{kyr}, लगभग एक अनुक्रिया काल है: हिमचादर के
दो-तिहाई रास्ता तय करने से पहले ही सूर्यातप पलट जाता है और उसे वापस खींच लेता है।
\textbf{22.} रैखिक मॉडल में समान आकार के प्रणोदनों के लिए बर्फ़ समान दर से बढ़ती और सिकुड़ती है:
बढ़ता और घटता प्रणोदन दर्पण-प्रतिबिंब अनुक्रियाएँ देते हैं। धीमे संचय के बाद तेज़
पतन के लिए ऐसी अनुक्रिया चाहिए जो अवस्था या परिवर्तन की दिशा पर निर्भर हो, यानी
एक अरैखिक अनुक्रिया।
\textbf{23.} हिमनदीय से अंतर्हिमनदीय तक कार्बन डाइऑक्साइड $3.93\log_2(266/194) = 1.8$ से $3.93\log_2(282/188) =
\qty{2.3}{W.m^{-2}}$ तक बढ़ी: पूरे ग्लोब पर और पूरे
वर्ष का प्रणोदन। अटलांटिक परिसंचरण के दुर्बल होने पर मुक्त होकर, वह उत्तरी
ग्रीष्म ऋतुओं के परिवर्तन को पूरे ग्रह तक ले जाती है, और सिकुड़ती बर्फ़ के एल्बिडो के साथ, उसे प्रवर्धित करती है।
\textbf{24.} जब अंतिम \omterm{def:b2:orbital-forcing:glacials}{अंतर्हिमनदीय काल} समाप्त हुआ, तब \ang{65} उ पर जून-अयनांत दैनिक-माध्य सूर्यातप
लगभग \textbf{\qty{439}{W.m^{-2}}} था, आज से \qty{39}{W.m^{-2}} कम और चक्र का अपना सबसे निम्न
मान।
\end{solution}
